\documentclass{abntex2}
\usepackage[utf8]{inputenc}
\usepackage{authblk}  % for better author formatting
\usepackage{amsmath}  % for mathematical typesetting
\usepackage{hyperref} % for clickable links like emails

% \titulo{Your Article Title}

% \autor[1]{Author One\thanks{Email: author.one@example.com}}
% \autor[2]{Author Two\thanks{Email: author.two@example.com}}

% \affil[1]{University of Brescia, Department of Quantitative Methods, Brescia, Italy}
% \affil[2]{École Polytechnique and GERAD, Montréal, Canada}
% \data{2024}

% \begin{document}
% \selectlanguage{brazil}
% \frenchspacing
% \maketitle
% \end{document}

% \titulo{Your Article Title}

% \autor{
% Equipe \abnTeX\thanks{``Recomenda-se que os dados de vinculação e
% endereço constem em nota, com sistema de chamada próprio, diferente do sistema 
% adotado para citações no texto.'' \url{http://www.abntex.net.br/}} 
% \\[0.5cm] 
% Lauro César Araujo\thanks{``Constar currículo sucinto de cada autor, com 
% vinculação corporativa e endereço de contato.''} }

% \begin{document}
% \maketitle
% \end{document}

% \titulo{Your Article Title}

% \author{Claudia Archetti\thanks{Some message}\affil{Institution 1} \and
%         Alain Hertz\thanks{Another message}\affil{Institution 2} \and
%         Maria Grazia Speranza\thanks{Additional info}\affil{Institution 1}}

% \begin{document}
% \maketitle

% \begin{abstract}
% This is the abstract of the document.
% \end{abstract}

% \section{Introduction}
% This is the introduction.

% \end{document}

\titulo{Your Article Title}

\autor{
    Lauro César Araujo\thanks{Email: lauro@example.com} \affil{Institution 1}   
    \\[0.5cm]
    Alain Hertz\thanks{Email: alain@example.com} \affil{Institution 2} 
    \\[0.5cm]
    Claudia Archetti\thanks{Email: claudia@example.com} \affil{Institution 1}
}

\begin{document}
\maketitle

\begin{abstract}
This is the abstract of the document.
\end{abstract}

\section{Introduction}
This is the introduction.

\end{document}
